\section{Discussion}
\label{sec:discussion}

\subsection{Coverage Trade-offs and the Value of Conditional Guarantees}

A recurring observation is that marginal CP baselines (DAPS, SNAPS) achieve
smaller APSS than FC-GNN on several datasets, but at the cost of
\emph{violated coverage} ($<\!90\%$ on most datasets).
This is not a mere numerical quirk: the calibration step of marginal CP
sets a global threshold that may be too permissive for rare attack communities
(giving small sets that miss the true class) and too conservative for common
benign communities.
Mondrian CP eliminates this bias by conditioning on community membership.
The larger APSS of FC-GNN relative to non-Mondrian models on some datasets
is therefore the \emph{expected cost} of providing genuine conditional
coverage — a cost that is mandatory in compliance-sensitive environments
(GDPR Article 22, NIS2 Article 21) where coverage failure constitutes a
material risk.

It is worth being precise about why a smaller prediction set is not
self-evidently better. Set size and coverage are not independent quantities
that a method may optimise separately: any method can drive APSS towards one by
raising its threshold, at the cost of dropping the true class more often. A set
size is therefore interpretable only once validity has been established, and
comparing APSS across methods that do not all attain the nominal level compares
quantities of different kinds. This is the reason we report coverage first and
efficiency second throughout Section~\ref{sec:experiments}, and the reason we
regard the coverage--efficiency frontier rather than either axis alone as the
correct object of comparison.

The practical consequence for a security operations centre is a difference in
failure mode rather than in average performance. A marginally calibrated
detector fails silently and non-uniformly: its aggregate coverage looks
acceptable while a specific attack family is systematically missed, and nothing
in the reported metrics reveals which family. A community-conditional detector
that returns a larger set fails visibly, by handing the analyst two or three
candidate classes instead of one. The second failure is far cheaper to absorb,
because it is legible at the point of triage.

\subsection{Scalability}

A detailed quantitative analysis of FC-GNN's scalability (training time, calibration time, and inference throughput) compared to baselines is provided in Section 4.8. While standard mini-batch training with neighbourhood sampling (GraphSAINT \citep{zeng20}) can scale the FMPL layer to $N \gg 10^4$ nodes, deploying FC-GNN in dynamic, high-throughput enterprise networks may still require incremental community detection algorithms to avoid re-computing Louvain communities from scratch during frequent graph updates.

The asymmetry between the two costs is worth noting, because it determines where
engineering effort should go. The fuzzy message-passing layer is a fixed
per-epoch cost that mini-batching addresses by standard means. The Louvain
partition is different in kind: it is comparatively cheap to compute once, but it
must be \emph{refreshed} whenever the graph changes materially, and the
conditional guarantee is defined relative to the partition in force at
calibration time. A stale partition does not merely degrade accuracy; it
invalidates the taxon on which the conditional coverage statement rests. The
binding constraint in a dynamic deployment is therefore partition maintenance
rather than model throughput, which is a further argument for hierarchical
partitioning schemes whose granularity can be adjusted incrementally.

\subsection{Exchangeability and Concept Drift}

The coverage theorem (Theorem~\ref{thm:coverage}) requires exchangeability
of Sugeno scores within each community, which holds exactly under the
fuzzy permutation-invariance condition (Definition~\ref{def:fpi}).
In practice, concept drift over time can violate this assumption:
as attack patterns evolve, nodes from the same community in the calibration
set may no longer be exchangeable with test nodes.
\citet{cp_drift} showed that CP pseudo-labels become inconsistent under
concept drift; our empirical results in Section 4.6 confirm this,
demonstrating that sliding-window recalibration is necessary to restore
conditional validity.
Recent work on Non-exchangeable Conformal Prediction \citep{ncpnet25} or
conformal change-point detection \citep{cpp_coad} offers theoretical
machinery to adjust $\Calpha$ dynamically without full recalibration,
which we leave for future integration.

We would stress that this is a limitation of the guarantee's scope rather than a
defect of the method, and that it is shared by every conformal procedure. The
theorem converts a distributional assumption into a finite-sample statement; it
does not create validity where the assumption fails. The appropriate operational
reading is that FC-GNN's guarantee holds \emph{between recalibrations}, with the
recalibration interval set by observed drift rather than fixed a priori. A
detector calibrated once and left running indefinitely does not retain its
guarantee, and reporting a coverage figure obtained under a random split — as is
customary in the CP-GNN literature — obscures this, because a random split
places calibration and test points in the same temporal regime by construction.

\subsection{Rule Complexity and Pruning}

The mean number of active rules per node (rule complexity) is very low
($0.002$--$0.023$) because the $\ell_1$ regularisation on firing strengths
$\|\bm{\tau}\|_1$ aggressively prunes inactive rules.
While this aids interpretability (fewer antecedents per explanation), it
may limit the expressiveness of the rule base for datasets with high
intra-class feature diversity (e.g., CIC-IDS-2017 with 15 classes).
Increasing $K$ from 16 to 32 or relaxing the $\ell_1$ penalty is expected
to improve Macro-F1 at a moderate cost in explanation simplicity.

The trade-off deserves to be stated as a design decision rather than a tuning
detail, because the two objectives are genuinely opposed. An explanation is
useful to an analyst in proportion to how few antecedents it cites; a classifier
is accurate on heterogeneous attack families in proportion to how many distinct
regions of feature space its rule base can represent. The sparsity penalty
selects a point on that curve, and no setting is optimal for both ends. We
report $K$ and the penalty weight explicitly so that the operating point is a
visible parameter of the method rather than an artefact of defaults.

\subsection{Limitations}

Despite these advantages, FC-GNN has several limitations:

\begin{enumerate}
\item \textbf{Point accuracy on CIC-IDS-2017.}
      FC-GNN achieves Macro-F1 $0.065$ on CIC-IDS-2017, reflecting the
      extreme class imbalance (${\sim}95\%$ benign).
      Class-weighted loss or SMOTE oversampling on the minority classes
      would improve this metric.
      We report the figure without mitigation so that all compared methods
      share one training protocol; a per-method tuned protocol would make the
      comparison across Table~\ref{tab:f1_main} incommensurable.
\item \textbf{No inductive evaluation on real PCAP data.}
      Our graph construction uses realistic synthetic statistics derived
      from the published feature distributions of each dataset.
      Direct evaluation on raw PCAP files requires a dedicated flow
      extraction pipeline (e.g., CICFlowMeter), which we leave for future
      work.
      This is the most consequential limitation of the study: real traffic
      differs from our construction in ways we do not model — heavy-tailed and
      non-Gaussian marginals, correlated and partially missing features, and
      topology shaped by routing and address translation rather than by a
      generative process. The experiments therefore support relative
      comparisons between methods under disclosed conditions, and should not be
      read as absolute performance claims on the named captures.
\item \textbf{Static $\lambda$ in Sugeno integral.}
      The $\lambda$ parameter is shared across all communities and
      trained globally.
      Per-community or per-class $\lambda$ values may better capture
      dataset-specific co-occurrence patterns among attack families. \rev{Given that attack categories exhibit vastly different structural characteristics (e.g., dense propagation in botnets versus sparse behavioral chains in malware infections), an adaptive, community-specific $\lambda$ would likely improve calibration efficiency and is an important avenue for future research.}
\item \rev{\textbf{Interpretability versus Causality.}
      While the fuzzy rules generated by FC-GNN provide empirical interpretability, they indicate structural \emph{correlation} rather than causal mechanisms. Under adversarial conditions, attackers may deliberately manipulate flow features to trigger misleading rules. Thus, the reliability of these explanations under adversarial feature manipulation requires further investigation.}
      \rev{We note that this concern is sharper for an interpretable model than
      for an opaque one: a rule base is an attack surface precisely because it
      is legible, since an adversary who can predict which rules fire can shape
      the justification presented to an analyst as well as the prediction
      itself.}
\item \textbf{Partition granularity.}
      Louvain is applied at its default resolution, and communities smaller
      than the minimum size fall back to the global quantile. Conditional
      coverage is consequently not uniform across the partition: the
      communities most likely to fall back are the small, rare ones, which are
      also those for which conditioning would be most valuable.
\end{enumerate}
